\documentclass{article}
\usepackage[utf8]{inputenc}
\usepackage[T1]{fontenc}
\usepackage{geometry}
\geometry{
    a4paper,
    left=2.5cm,
    right=2.5cm,
    top=2.5cm,
    bottom=2.5cm
}
\usepackage{graphicx}
\usepackage{float}
\usepackage{hyperref}
\usepackage[english]{babel}
\usepackage{amsmath}
\usepackage{amssymb}
\usepackage{array}
\usepackage{booktabs}
\usepackage{caption}
\usepackage{subcaption}
\usepackage{enumitem}
\usepackage{fancyhdr}
\usepackage{titlesec}
\usepackage{url}
\usepackage{textcomp}
\usepackage{xcolor}
\usepackage{listings}
\usepackage{multirow}
\usepackage{longtable}
\usepackage{tabularx}
\usepackage{times}

% Set up page style
\pagestyle{fancy}
\fancyhf{}
\rfoot{\thepage}
\renewcommand{\headrulewidth}{0pt}

% Title formatting
\titleformat{\chapter}[display]{\normalfont\bfseries}{\chaptertitle\thechapter}{0pt}{\Huge}
\titlespacing*{\chapter}{0pt}{50pt}{40pt}

% Document info
\title{Template Document}
\author{Academic Support Agent}
\date{\today}

\begin{document}
\maketitle

\section*{La princesse de cl ve}
	(1678) Une esth tique classique
\begin{enumerate}
   \item Une incarnation de la perfection
   \begin{itemize}
     \item Le portrait physique
     \item Fine fleur de l'aristocratie
     \item L'importance du portrait moral
   \end{itemize}
   \item Une  ducation irr prochable
   \begin{itemize}
       \item  loge de Mme de Chartres charg e de l' ducation de sa fille
       \item Pr paration   la vie de cour et aux risques qu'elle comporte
       \item Une vision exigeante de l'amour et de l'honn te femme
    \end{itemize}
\end{enumerate}
\noindent\hrulefill

\section*{La princesse de cl ve}
	(1678) Une esth tique classique
\begin{enumerate}
   \item Une incarnation de la perfection
   \begin{itemize}
     \item Le portrait physique
     \item Fine fleur de l'aristocratie
     \item L'importance du portrait moral
   \end{itemize}
   \item Une  ducation irr prochable
   \begin{itemize}
       \item  loge de Mme de Chartres charg e de l' ducation de sa fille
       \item Pr paration   la vie de cour et aux risques qu'elle comporte
       \item Une vision exigeante de l'amour et de l'honn te femme
    \end{itemize}
\end{enumerate}
\noindent\hrulefill

\section*{La princesse de cl ve}
	(1678) Une esth tique classique
\begin{enumerate}
   \item Une incarnation de la perfection
   \begin{itemize}
     \item Le portrait physique
     \item Fine fleur de l'aristocratie
     \item L'importance du portrait moral
   \end{itemize}
   \item Une  ducation irr prochable
   \begin{itemize}
       \item  loge de Mme de Chartres charg e de l' ducation de sa fille
       \item Pr paration   la vie de cour et aux risques qu'elle comporte
       \item Une vision exigeante de l'amour et de l'honn te femme
    \end{itemize}
\end{enumerate}

\section*{La princesse de cl ve}
	(1678) Une esth tique classique
\begin{enumerate}
   \item Une incarnation de la perfection
   \begin{itemize}
     \item Le portrait physique
     \item Fine fleur de l'aristocratie
     \item L'importance du portrait moral
   \end{itemize}
   \item Une  ducation irr prochable
   \begin{itemize}
       \item  loge de Mme de Chartres charg e de l' ducation de sa fille
       \item Pr paration   la vie de cour et aux risques qu'elle comporte
       \item Une vision exigeante de l'amour et de l'honn te femme
    \end{itemize}
\end{enumerate}
\noindent\hrulefill

\section*{La princesse de cl ve}
	(1678) Une esth tique classique
\begin{enumerate}
   \item Une incarnation de la perfection
   \begin{itemize}
     \item Le portrait physique
     \item Fine fleur de l'aristocratie
     \item L'importance du portrait moral
   \end{itemize}
   \item Une  ducation irr prochable
   \begin{itemize}
       \item  loge de Mme de Chartres charg e de l' ducation de sa fille
       \item Pr paration   la vie de cour et aux risques qu'elle comporte
       \item Une vision exigeante de l'amour et de l'honn te femme
    \end{itemize}
\end{enumerate}
\noindent\hrulefill

\section*{La princesse de cl ve}
	(1678) Une esth tique classique
\begin{enumerate}
   \item Une incarnation de la perfection
   \begin{itemize}
     \item Le portrait physique
     \item Fine fleur de l'aristocratie
     \item L'importance du portrait moral
   \end{itemize}
   \item Une  ducation irr prochable
   \begin{itemize}
       \item  loge de Mme de Chartres charg e de l' ducation de sa fille
       \item Pr paration   la vie de cour et aux risques qu'elle comporte
       \item Une vision exigeante de l'amour et de l'honn te femme
    \end{itemize}
\end{enumerate}

\section*{La princesse de cl ve}
	(1678) Une esth tique classique
\begin{enumerate}
   \item Une incarnation de la perfection
   \begin{itemize}
     \item Le portrait physique
     \item Fine fleur de l'aristocratie
     \item L'importance du portrait moral
   \end{itemize}
   \item Une  ducation irr prochable
   \begin{itemize}
       \item  loge de Mme de Chartres charg e de l' ducation de sa fille
       \item Pr paration   la vie de cour et aux risques qu'elle comporte
       \item Une vision exigeante de l'amour et de l'honn te femme
    \end{itemize}
\end{enumerate}
\noindent\hrulefill

\section*{La princesse de cl ve}
	(1678) Une esth tique classique
\begin{enumerate}
   \item Une incarnation de la perfection
   \begin{itemize}
     \item Le portrait physique
     \item Fine fleur de l'aristocratie
     \item L'importance du portrait moral
   \end{itemize}
   \item Une  ducation irr prochable
   \begin{itemize}
       \item  loge de Mme de Chartres charg e de l' ducation de sa fille
       \item Pr paration   la vie de cour et aux risques qu'elle comporte
       \item Une vision exigeante de l'amour et de l'honn te femme
    \end{itemize}
\end{enumerate}
\noindent\hrulefill

\section*{La princesse de cl ve}
	(1678) Une esth tique classique
\begin{enumerate}
   \item Une incarnation de la perfection
   \begin{itemize}
     \item Le portrait physique
     \item Fine fleur de l'aristocratie
     \item L'importance du portrait moral
   \end{itemize}
   \item Une  ducation irr prochable
   \begin{itemize}
       \item  loge de Mme de Chartres charg e de l' ducation de sa fille
       \item Pr paration   la vie de cour et aux risques qu'elle comporte
       \item Une vision exigeante de l'amour et de l'honn te femme
    \end{itemize}
\end{enumerate}
\end{document}